
\documentclass[a4paper,12pt]{article}

\usepackage[T2A]{fontenc}
\usepackage[utf8]{inputenc}
\usepackage[russian]{babel}

\usepackage{amsmath,amssymb,mathtools}
\usepackage{helvet}
\renewcommand{\familydefault}{\sfdefault}
\usepackage{icomma}
\usepackage[left=18mm,right=18mm,top=18mm,bottom=20mm]{geometry}
\usepackage{microtype}
\usepackage{tabularx,booktabs,array}
\usepackage{enumitem}
\usepackage{fancyhdr}
\usepackage{setspace}
\usepackage{xcolor}
\usepackage{hyperref}

\hypersetup{hidelinks}

\definecolor{accent}{HTML}{008080}
\definecolor{darkaccent}{HTML}{004D4D}
\definecolor{textgray}{HTML}{333333}
\definecolor{lightline}{HTML}{D9EAEA}

\onehalfspacing
\setlength{\parindent}{0pt}
\setlength{\parskip}{6pt}

\pagestyle{fancy}
\fancyhf{}
\lhead{\textcolor{darkaccent}{Тригонометрия}}
\rhead{\textcolor{darkaccent}{Тимофей}}
\cfoot{\thepage}

\newcommand{\sectionline}{%
\vspace{1mm}
\noindent\textcolor{lightline}{\rule{\linewidth}{0.8pt}}
\vspace{1mm}
}

\newcommand{\answer}[1]{\textbf{Ответ: }#1}
\newcommand{\important}[1]{\textcolor{darkaccent}{\textbf{#1}}}

\begin{document}

\begin{titlepage}
\thispagestyle{empty}

\vspace*{18mm}

{\Huge\bfseries\textcolor{darkaccent}{Чек-лист}}

\vspace{8mm}

{\Large\bfseries Тригонометрия: сумма синуса и косинуса}

\vspace{4mm}

{\large Формулы приведения, вспомогательный угол и замена через половинный угол}

\vspace{18mm}

\textbf{Ученик:} Тимофей Васильченко

\vspace{7mm}

\textbf{Цель:} научиться решать уравнения вида
\[
a\sin\alpha+b\cos\alpha=c
\]
несколькими способами и понимать, когда уравнение не имеет решений.

\vfill

{\large Лёвин Артём Александрович}

\vspace{4mm}

{\large \today}

\end{titlepage}

\section*{1. Главная идея занятия}

На занятии разбирались уравнения и преобразования, где в одном выражении встречаются синус и косинус:
\[
\sin x+\cos x,\qquad \sin x-\cos x,\qquad a\sin\alpha+b\cos\alpha=c.
\]

Главная мысль: если выражение не подходит под готовую формулу, его нужно \important{подогнать под формулу}. Для этого используются:

\begin{enumerate}[label=\textbf{\arabic*.}, leftmargin=8mm]
\item формулы приведения;
\item формулы суммы и разности тригонометрических функций;
\item метод вспомогательного угла;
\item переход к половинному углу.
\end{enumerate}

\section*{2. Формулы приведения}

Чтобы заменить синус на косинус или косинус на синус, используем:

\[
\sin x=\cos\left(\frac{\pi}{2}-x\right),
\]
\[
\cos x=\sin\left(\frac{\pi}{2}-x\right).
\]

\sectionline

\textbf{Зачем это нужно?}

Например, выражение
\[
\sin x-\cos x
\]
нельзя сразу свернуть по формуле разности синусов, потому что здесь синус и косинус.

Но можно заменить:
\[
\cos x=\sin\left(\frac{\pi}{2}-x\right).
\]

Тогда:
\[
\sin x-\cos x
=
\sin x-\sin\left(\frac{\pi}{2}-x\right).
\]

Теперь уже можно применять формулу разности синусов.

\section*{3. Формулы суммы и разности}

\[
\sin A+\sin B
=
2\sin\frac{A+B}{2}\cos\frac{A-B}{2}.
\]

\[
\sin A-\sin B
=
2\cos\frac{A+B}{2}\sin\frac{A-B}{2}.
\]

\[
\cos A+\cos B
=
2\cos\frac{A+B}{2}\cos\frac{A-B}{2}.
\]

\[
\cos A-\cos B
=
-2\sin\frac{A+B}{2}\sin\frac{A-B}{2}.
\]

\sectionline

\textbf{Пример. Свернём выражение \(\sin x-\cos x\).}

\[
\sin x-\cos x
=
\sin x-\sin\left(\frac{\pi}{2}-x\right).
\]

Применяем формулу:
\[
\sin A-\sin B
=
2\cos\frac{A+B}{2}\sin\frac{A-B}{2}.
\]

Получаем:
\[
\sin x-\sin\left(\frac{\pi}{2}-x\right)
=
2\cos\frac{x+\frac{\pi}{2}-x}{2}
\sin\frac{x-\left(\frac{\pi}{2}-x\right)}{2}.
\]

Упрощаем:
\[
=
2\cos\frac{\pi}{4}
\sin\left(x-\frac{\pi}{4}\right).
\]

Так как
\[
\cos\frac{\pi}{4}=\frac{\sqrt2}{2},
\]
получаем:
\[
\sin x-\cos x
=
\sqrt2\sin\left(x-\frac{\pi}{4}\right).
\]

\[
\boxed{\sin x-\cos x=\sqrt2\sin\left(x-\frac{\pi}{4}\right)}
\]

\section*{4. Частный случай: \(\sin x+\cos x\)}

Очень полезная формула:
\[
\sin x+\cos x
=
\sqrt2\sin\left(x+\frac{\pi}{4}\right).
\]

Также можно записать:
\[
\sin x+\cos x
=
\sqrt2\cos\left(x-\frac{\pi}{4}\right).
\]

\sectionline

\textbf{Как это получить?}

\[
\sin x+\cos x
=
\sin x+\sin\left(\frac{\pi}{2}-x\right).
\]

По формуле суммы синусов:
\[
\sin x+\sin\left(\frac{\pi}{2}-x\right)
=
2\sin\frac{x+\frac{\pi}{2}-x}{2}
\cos\frac{x-\left(\frac{\pi}{2}-x\right)}{2}.
\]

\[
=
2\sin\frac{\pi}{4}\cos\left(x-\frac{\pi}{4}\right).
\]

Так как
\[
\sin\frac{\pi}{4}=\frac{\sqrt2}{2},
\]
получаем:
\[
\sin x+\cos x
=
\sqrt2\cos\left(x-\frac{\pi}{4}\right).
\]

\section*{5. Уравнение \(\sin x+\cos x=1\)}

Используем формулу:
\[
\sin x+\cos x
=
\sqrt2\cos\left(x-\frac{\pi}{4}\right).
\]

Тогда:
\[
\sqrt2\cos\left(x-\frac{\pi}{4}\right)=1.
\]

Делим на \(\sqrt2\):
\[
\cos\left(x-\frac{\pi}{4}\right)=\frac{1}{\sqrt2}.
\]

Так как
\[
\frac{1}{\sqrt2}=\frac{\sqrt2}{2},
\]
получаем:
\[
\cos\left(x-\frac{\pi}{4}\right)=\frac{\sqrt2}{2}.
\]

Отсюда:
\[
x-\frac{\pi}{4}=\pm\frac{\pi}{4}+2\pi n,\qquad n\in\mathbb Z.
\]

Значит:
\[
x=\frac{\pi}{4}\pm\frac{\pi}{4}+2\pi n.
\]

Получаем две серии:
\[
x=2\pi n,
\]
\[
x=\frac{\pi}{2}+2\pi n,
\qquad n\in\mathbb Z.
\]

\[
\answer{\displaystyle x=2\pi n \quad \text{или} \quad x=\frac{\pi}{2}+2\pi n,\quad n\in\mathbb Z.}
\]

\section*{6. Метод вспомогательного угла}

Рассмотрим общий вид:
\[
a\sin\alpha+b\cos\alpha=c.
\]

Вводим число:
\[
R=\sqrt{a^2+b^2}.
\]

Делим всё уравнение на \(R\):
\[
\frac{a}{R}\sin\alpha+\frac{b}{R}\cos\alpha=\frac{c}{R}.
\]

Теперь вводим угол \(\varphi\) так, чтобы:
\[
\cos\varphi=\frac{a}{R},
\]
\[
\sin\varphi=\frac{b}{R}.
\]

Тогда:
\[
\frac{a}{R}\sin\alpha+\frac{b}{R}\cos\alpha
=
\cos\varphi\sin\alpha+\sin\varphi\cos\alpha.
\]

По формуле синуса суммы:
\[
\sin(\alpha+\varphi)
=
\sin\alpha\cos\varphi+\cos\alpha\sin\varphi.
\]

Значит:
\[
\cos\varphi\sin\alpha+\sin\varphi\cos\alpha
=
\sin(\alpha+\varphi).
\]

И уравнение превращается в:
\[
\sin(\alpha+\varphi)=\frac{c}{R}.
\]

\section*{7. Когда решений нет}

Так как синус любого угла лежит от \(-1\) до \(1\):
\[
-1\le\sin(\alpha+\varphi)\le1,
\]
уравнение
\[
\sin(\alpha+\varphi)=\frac{c}{R}
\]
имеет решения только тогда, когда:
\[
\left|\frac{c}{R}\right|\le1.
\]

Если:
\[
\left|\frac{c}{R}\right|>1,
\]
то решений нет.

\sectionline

\textbf{Пример.}

\[
\sin\alpha+2\cos\alpha=3.
\]

Здесь:
\[
a=1,\qquad b=2,\qquad c=3.
\]

\[
R=\sqrt{1^2+2^2}=\sqrt5.
\]

Правая часть после деления:
\[
\frac{c}{R}=\frac3{\sqrt5}.
\]

Так как:
\[
\frac3{\sqrt5}>1,
\]
уравнение решений не имеет.

\[
\answer{\displaystyle \varnothing.}
\]

\section*{8. Пример метода вспомогательного угла}

Решим уравнение:
\[
2\sin\alpha+3\cos\alpha=1.
\]

Здесь:
\[
a=2,\qquad b=3,\qquad c=1.
\]

\[
R=\sqrt{2^2+3^2}=\sqrt{13}.
\]

Делим уравнение на \(\sqrt{13}\):
\[
\frac{2}{\sqrt{13}}\sin\alpha+
\frac{3}{\sqrt{13}}\cos\alpha
=
\frac{1}{\sqrt{13}}.
\]

Вводим угол \(\varphi\):
\[
\cos\varphi=\frac{2}{\sqrt{13}},
\]
\[
\sin\varphi=\frac{3}{\sqrt{13}}.
\]

Тогда:
\[
\sin(\alpha+\varphi)=\frac{1}{\sqrt{13}}.
\]

Поскольку:
\[
\left|\frac{1}{\sqrt{13}}\right|<1,
\]
решения есть.

Общее решение:
\[
\alpha+\varphi=\arcsin\frac{1}{\sqrt{13}}+2\pi n,
\]
или
\[
\alpha+\varphi=\pi-\arcsin\frac{1}{\sqrt{13}}+2\pi n,
\qquad n\in\mathbb Z.
\]

Значит:
\[
\alpha=\arcsin\frac{1}{\sqrt{13}}-\varphi+2\pi n,
\]
или
\[
\alpha=\pi-\arcsin\frac{1}{\sqrt{13}}-\varphi+2\pi n.
\]

Так как:
\[
\varphi=\arccos\frac{2}{\sqrt{13}},
\]
получаем:
\[
\answer{
\displaystyle
\alpha=\arcsin\frac{1}{\sqrt{13}}-\arccos\frac{2}{\sqrt{13}}+2\pi n
}
\]
или
\[
\answer{
\displaystyle
\alpha=\pi-\arcsin\frac{1}{\sqrt{13}}-\arccos\frac{2}{\sqrt{13}}+2\pi n,\quad n\in\mathbb Z.
}
\]

\section*{9. Второй способ: переход к половинному углу}

Уравнение вида
\[
a\sin\alpha+b\cos\alpha=c
\]
можно решать через формулы:
\[
\sin\alpha=2\sin\frac{\alpha}{2}\cos\frac{\alpha}{2},
\]
\[
\cos\alpha=\cos^2\frac{\alpha}{2}-\sin^2\frac{\alpha}{2},
\]
\[
1=\sin^2\frac{\alpha}{2}+\cos^2\frac{\alpha}{2}.
\]

После этого часто получается однородное уравнение относительно:
\[
\sin\frac{\alpha}{2}
\quad \text{и} \quad
\cos\frac{\alpha}{2}.
\]

Затем можно делить на
\[
\cos^2\frac{\alpha}{2},
\]
но только после отдельной проверки случая:
\[
\cos\frac{\alpha}{2}=0.
\]

\section*{10. Важное замечание про ОДЗ}

Исходное уравнение
\[
2\sin\alpha+3\cos\alpha=1
\]
определено при всех \(\alpha\).

Поэтому у него:
\[
\alpha\in\mathbb R.
\]

Если в процессе решения мы делим на
\[
\cos^2\frac{\alpha}{2},
\]
то условие
\[
\cos\frac{\alpha}{2}\ne0
\]
не является ОДЗ исходного уравнения.

Это техническое условие для конкретного преобразования.

Правильно писать так:

\[
\cos\frac{\alpha}{2}=0
\]
проверим отдельно.

Если этот случай не даёт решений, тогда можно продолжать решение при условии:
\[
\cos\frac{\alpha}{2}\ne0.
\]

\section*{11. Пример через половинный угол}

Решим:
\[
2\sin\alpha+3\cos\alpha=1.
\]

Подставляем:
\[
\sin\alpha=2\sin\frac{\alpha}{2}\cos\frac{\alpha}{2},
\]
\[
\cos\alpha=\cos^2\frac{\alpha}{2}-\sin^2\frac{\alpha}{2},
\]
\[
1=\sin^2\frac{\alpha}{2}+\cos^2\frac{\alpha}{2}.
\]

Получаем:
\[
2\cdot2\sin\frac{\alpha}{2}\cos\frac{\alpha}{2}
+
3\left(
\cos^2\frac{\alpha}{2}
-
\sin^2\frac{\alpha}{2}
\right)
=
\sin^2\frac{\alpha}{2}+\cos^2\frac{\alpha}{2}.
\]

Раскрываем скобки:
\[
4\sin\frac{\alpha}{2}\cos\frac{\alpha}{2}
+
3\cos^2\frac{\alpha}{2}
-
3\sin^2\frac{\alpha}{2}
=
\sin^2\frac{\alpha}{2}+\cos^2\frac{\alpha}{2}.
\]

Переносим всё влево:
\[
4\sin\frac{\alpha}{2}\cos\frac{\alpha}{2}
+
2\cos^2\frac{\alpha}{2}
-
4\sin^2\frac{\alpha}{2}
=0.
\]

Делим на \(2\):
\[
2\sin\frac{\alpha}{2}\cos\frac{\alpha}{2}
+
\cos^2\frac{\alpha}{2}
-
2\sin^2\frac{\alpha}{2}
=0.
\]

Проверяем случай:
\[
\cos\frac{\alpha}{2}=0.
\]

Тогда:
\[
\frac{\alpha}{2}=\frac{\pi}{2}+\pi n,
\]
\[
\alpha=\pi+2\pi n.
\]

Подставим в исходное уравнение:
\[
2\sin(\pi+2\pi n)+3\cos(\pi+2\pi n)=0+3\cdot(-1)=-3\ne1.
\]

Значит, этот случай решений не даёт.

Теперь можно делить на
\[
\cos^2\frac{\alpha}{2}.
\]

Получаем:
\[
2\tg\frac{\alpha}{2}
+
1
-
2\tg^2\frac{\alpha}{2}
=0.
\]

Пусть:
\[
t=\tg\frac{\alpha}{2}.
\]

Тогда:
\[
2t+1-2t^2=0.
\]

Умножим на \(-1\):
\[
2t^2-2t-1=0.
\]

Решаем:
\[
D=(-2)^2-4\cdot2\cdot(-1)=4+8=12.
\]

\[
t_{1,2}=\frac{2\pm\sqrt{12}}{4}
=
\frac{2\pm2\sqrt3}{4}
=
\frac{1\pm\sqrt3}{2}.
\]

Возвращаемся к \(\alpha\):
\[
\tg\frac{\alpha}{2}=\frac{1+\sqrt3}{2}
\]
или
\[
\tg\frac{\alpha}{2}=\frac{1-\sqrt3}{2}.
\]

Отсюда:
\[
\frac{\alpha}{2}
=
\arctg\frac{1+\sqrt3}{2}+\pi n,
\]
или
\[
\frac{\alpha}{2}
=
\arctg\frac{1-\sqrt3}{2}+\pi n.
\]

Значит:
\[
\alpha
=
2\arctg\frac{1+\sqrt3}{2}+2\pi n,
\]
или
\[
\alpha
=
2\arctg\frac{1-\sqrt3}{2}+2\pi n,
\qquad n\in\mathbb Z.
\]

\[
\answer{
\displaystyle
\alpha
=
2\arctg\frac{1+\sqrt3}{2}+2\pi n
}
\]
или
\[
\answer{
\displaystyle
\alpha
=
2\arctg\frac{1-\sqrt3}{2}+2\pi n,\quad n\in\mathbb Z.
}
\]

\section*{12. Какой метод выбирать}

\begin{center}
\renewcommand{\arraystretch}{1.45}
\begin{tabularx}{\linewidth}{p{0.32\linewidth}X}
\toprule
\textbf{Ситуация} & \textbf{Удобный метод} \\
\midrule
\(\sin x+\cos x\), \(\sin x-\cos x\) & Формулы приведения и формулы суммы/разности \\
\midrule
\(a\sin\alpha+b\cos\alpha=c\) с любыми коэффициентами & Метод вспомогательного угла \\
\midrule
Коэффициенты простые, хочется получить квадратное уравнение & Переход к половинному углу \\
\midrule
Правая часть после деления на \(R\) больше \(1\) по модулю & Сразу пишем: решений нет \\
\bottomrule
\end{tabularx}
\end{center}

\section*{13. Типичные ошибки}

\subsection*{Ошибка 1. Смешать ОДЗ и техническое условие}

Нельзя писать:
\[
\cos\frac{\alpha}{2}\ne0
\]
как ОДЗ исходного уравнения
\[
2\sin\alpha+3\cos\alpha=1.
\]

Исходное уравнение определено при всех \(\alpha\).

Правильно: сначала отдельно проверить случай
\[
\cos\frac{\alpha}{2}=0,
\]
а потом делить на
\[
\cos^2\frac{\alpha}{2}.
\]

\subsection*{Ошибка 2. Забыть проверить отсутствие решений}

Для уравнения
\[
a\sin\alpha+b\cos\alpha=c
\]
нужно считать:
\[
R=\sqrt{a^2+b^2}.
\]

Если:
\[
\left|\frac{c}{R}\right|>1,
\]
то решений нет.

\subsection*{Ошибка 3. Неправильно ввести угол \(\varphi\)}

Если:
\[
a\sin\alpha+b\cos\alpha=c,
\]
то после деления на \(R\):
\[
\frac{a}{R}\sin\alpha+\frac{b}{R}\cos\alpha=\frac{c}{R}.
\]

Чтобы получить синус суммы, нужно брать:
\[
\cos\varphi=\frac{a}{R},
\]
\[
\sin\varphi=\frac{b}{R}.
\]

Тогда:
\[
\frac{a}{R}\sin\alpha+\frac{b}{R}\cos\alpha
=
\sin(\alpha+\varphi).
\]

\subsection*{Ошибка 4. Потерять второй корень уравнения с синусом}

Если:
\[
\sin y=m,
\]
то общее решение:
\[
y=\arcsin m+2\pi n
\]
или
\[
y=\pi-\arcsin m+2\pi n,
\qquad n\in\mathbb Z.
\]

Нужно записывать обе серии.

\subsection*{Ошибка 5. Делить на выражение, которое может быть равно нулю}

Перед делением на
\[
\cos\frac{\alpha}{2}
\quad \text{или} \quad
\cos^2\frac{\alpha}{2}
\]
нужно отдельно проверить случай:
\[
\cos\frac{\alpha}{2}=0.
\]

\newpage

\section*{Тренировочный блок}

\subsection*{Задание 1}

Сверните выражение в одну тригонометрическую функцию:
\[
\sin x+\cos x.
\]

\subsection*{Задание 2}

Сверните выражение в одну тригонометрическую функцию:
\[
\sin x-\cos x.
\]

\subsection*{Задание 3}

Решите уравнение:
\[
\sin x+\cos x=1.
\]

\subsection*{Задание 4}

Решите уравнение:
\[
\sin\alpha+2\cos\alpha=3.
\]

\subsection*{Задание 5}

Решите уравнение методом вспомогательного угла:
\[
\sin\alpha+2\cos\alpha=1.
\]

\subsection*{Задание 6}

Решите уравнение методом вспомогательного угла:
\[
2\sin\alpha+3\cos\alpha=1.
\]

\subsection*{Задание 7}

Решите уравнение через половинный угол:
\[
\sin\alpha+2\cos\alpha=1.
\]

\subsection*{Задание 8}

Решите уравнение:
\[
3\sin x+4\cos x=5.
\]

\subsection*{Задание 9}

Решите уравнение:
\[
3\sin x+4\cos x=6.
\]

\subsection*{Задание 10}

Проверьте, имеет ли решения уравнение:
\[
5\sin x-12\cos x=13.
\]

\newpage

\section*{Ответы}

\begin{center}
\renewcommand{\arraystretch}{1.5}
\begin{tabularx}{\linewidth}{>{\centering\arraybackslash}p{14mm}X}
\toprule
\textbf{№} & \textbf{Ответ} \\
\midrule
1 & \(\displaystyle \sin x+\cos x=\sqrt2\sin\left(x+\frac{\pi}{4}\right)\) \\
\midrule
2 & \(\displaystyle \sin x-\cos x=\sqrt2\sin\left(x-\frac{\pi}{4}\right)\) \\
\midrule
3 & \(\displaystyle x=2\pi n\) или \(\displaystyle x=\frac{\pi}{2}+2\pi n,\quad n\in\mathbb Z\) \\
\midrule
4 & \(\displaystyle \varnothing\) \\
\midrule
5 &
\(\displaystyle
\alpha=\arcsin\frac1{\sqrt5}-\arccos\frac1{\sqrt5}+2\pi n
\)
или

\(\displaystyle
\alpha=\pi-\arcsin\frac1{\sqrt5}-\arccos\frac1{\sqrt5}+2\pi n,\quad n\in\mathbb Z
\)
\\
\midrule
6 &
\(\displaystyle
\alpha=\arcsin\frac1{\sqrt{13}}-\arccos\frac2{\sqrt{13}}+2\pi n
\)
или

\(\displaystyle
\alpha=\pi-\arcsin\frac1{\sqrt{13}}-\arccos\frac2{\sqrt{13}}+2\pi n,\quad n\in\mathbb Z
\)
\\
\midrule
7 &
\(\displaystyle
\alpha=2\arctg 1+2\pi n
\)
или

\(\displaystyle
\alpha=2\arctg\left(-\frac13\right)+2\pi n,\quad n\in\mathbb Z
\)
\\
\midrule
8 &
\(\displaystyle
x=\frac{\pi}{2}-\varphi+2\pi n,\quad n\in\mathbb Z,
\)
где
\(\displaystyle
\cos\varphi=\frac35,\quad \sin\varphi=\frac45
\)
\\
\midrule
9 & \(\displaystyle \varnothing\) \\
\midrule
10 & \(\displaystyle x=-\arccos\frac5{13}+2\pi n,\quad n\in\mathbb Z\) \\
\bottomrule
\end{tabularx}
\end{center}

\section*{Короткая памятка}

\begin{enumerate}[label=\textbf{\arabic*.}, leftmargin=8mm]
\item \(\sin x=\cos\left(\frac{\pi}{2}-x\right)\), \(\cos x=\sin\left(\frac{\pi}{2}-x\right)\).
\item \(\sin x+\cos x=\sqrt2\sin\left(x+\frac{\pi}{4}\right)\).
\item \(\sin x-\cos x=\sqrt2\sin\left(x-\frac{\pi}{4}\right)\).
\item Для \(a\sin\alpha+b\cos\alpha=c\) считай \(R=\sqrt{a^2+b^2}\).
\item Если \(\left|\frac{c}{R}\right|>1\), решений нет.
\item Если \(\left|\frac{c}{R}\right|\le1\), вводи угол \(\varphi\):
\[
\cos\varphi=\frac{a}{R},\qquad \sin\varphi=\frac{b}{R}.
\]
\item Тогда:
\[
a\sin\alpha+b\cos\alpha=c
\quad \Longleftrightarrow \quad
\sin(\alpha+\varphi)=\frac{c}{R}.
\]
\item При делении на \(\cos\frac{\alpha}{2}\) или \(\cos^2\frac{\alpha}{2}\) сначала отдельно проверь случай \(\cos\frac{\alpha}{2}=0\).
\item Не называй техническое условие деления ОДЗ исходного уравнения.
\item Для \(\sin y=m\) всегда записывай две серии решений.
\end{enumerate}

\end{document}

